\begin{frame}{$V^\pi$ 와 $Q^\pi$: 동전의 양면, 그리고 2$\times$2}
  \kb{행동가치함수} $Q^\pi(s,a)$ --- ``상태 $s$ 에서 \textbf{행동
  $a$ 를 강제 한 번 취하고}, 그 다음부터 $\pi$ 를 따른다''의
  가치 --- 는 $V^\pi$ 와 같은 방식으로 재귀되고, 다리
  $Q^\pi(s,\pi(s)) = V^\pi(s)$ 로 서로를 정의하는 \kb{동전의
  양면}:
  \[
  Q^\pi(s,a) = R(s,a) + \gamma \sum_{s'} P(s'|s,a)\, V^\pi(s')
  \]
  \begin{columns}[T]
    \begin{column}{0.44\textwidth}
      \centering
      \includegraphics[width=0.95\linewidth]{ch04_1_bellman_equations.png}
      \vskip0.3em
      {\footnotesize 네 개의 벨만방정식: \textbf{기대 vs
      최적} × \textbf{$V$ vs $Q$} 로 2$\times$2.}
    \end{column}
    \begin{column}{0.56\textwidth}
      \small
      \begin{tabular}{@{}lcc@{}}
        \toprule
         & \kb{기대}($\pi$ 를 따른다) & \kb{최적}($\max$) \\
        \midrule
        \textbf{$V$} & $V^\pi$ (4.1절 = 오늘) & $V^{*}$ (4.2--4.3절) \\
        \textbf{$Q$} & $Q^\pi$ & $Q^{*}$ (Week 9 DQN) \\
        \bottomrule
      \end{tabular}
      \vskip0.4em
      {\footnotesize 이번 절의 모든 코드와 수식이 이 다리(V-방정식
      $\leftrightarrow$ Q-방정식 사이)를 통해 오간다.}
    \end{column}
  \end{columns}
\end{frame}
